\section{Sonuçların Değerlendirilmesi}

Bu çalışmada, Hakkari bölgesi için günlük yıldırım tahmini amacıyla dört farklı makine öğrenmesi algoritması karşılaştırılmıştır. Elde edilen sonuçlar, makine öğrenmesi yöntemlerinin yıldırım tahmini için başarılı bir şekilde kullanılabileceğini göstermektedir.

\subsection{Model Performansları}

Karşılaştırılan modeller arasında:

\begin{itemize}
    \item En yüksek F1-Score değeri [model adı] tarafından elde edilmiştir.
    \item Tree-based modeller (Random Forest, XGBoost, LightGBM) genel olarak Lojistik Regresyon'dan daha iyi performans göstermiştir.
    \item Sınıf dengesizliği, modellerin precision ve recall değerleri arasında trade-off oluşturmaktadır.
\end{itemize}

\subsection{Özellik Önemleri}

Özellik önem analizi, yıldırım tahmininde en etkili faktörlerin:

\begin{enumerate}
    \item Önceki günlerin yıldırım durumu
    \item Mevsimsel özellikler
    \item Coğrafi özellikler (rakım)
\end{enumerate}

olduğunu ortaya koymuştur. Bu bulgular, yıldırım olaylarının belirli meteorolojik koşullar altında kümeleme eğilimi gösterdiğini desteklemektedir.

\section{Çalışmanın Katkıları}

Bu çalışmanın başlıca katkıları:

\begin{enumerate}
    \item Hakkari bölgesi için ilk kapsamlı yıldırım tahmin çalışması
    \item 10 yıllık uzun dönemli veri analizi
    \item Farklı makine öğrenmesi algoritmalarının karşılaştırmalı değerlendirmesi
    \item Pratik uygulamalar için temel oluşturacak özellik mühendisliği yaklaşımları
\end{enumerate}

\section{Sınırlılıklar}

Bu çalışmanın bazı sınırlılıkları bulunmaktadır:

\begin{itemize}
    \item Meteorolojik veriler günlük frekansta olup saatlik tahmin yapmaya olanak vermemektedir.
    \item Bazı meteorolojik değişkenlerde (nem, bulutluluk) eksik veriler mevcuttur.
    \item Model sadece Hakkari bölgesi için eğitilmiş olup diğer bölgelere genellenebilirliği test edilmemiştir.
\end{itemize}

\section{Gelecek Çalışmalar}

Bu çalışma, aşağıdaki yönlerde genişletilebilir:

\begin{enumerate}
    \item \textbf{Saatlik Tahmin:} Daha yüksek çözünürlüklü verilerle saatlik tahmin modelleri
    \item \textbf{Radar Verileri:} Radar yansıtıcılık verilerinin modele eklenmesi
    \item \textbf{Derin Öğrenme:} LSTM ve CNN gibi derin öğrenme mimarilerinin denenmesi
    \item \textbf{Bölgesel Genelleme:} Modelin Türkiye'nin diğer bölgelerine uyarlanması
    \item \textbf{Operasyonel Sistem:} Gerçek zamanlı tahmin sistemi geliştirme
\end{enumerate}

\section{Sonuç}

Bu çalışmada, Hakkari bölgesi için meteorolojik veriler kullanılarak yıldırım tahmini yapılmıştır. Makine öğrenmesi algoritmalarının yıldırım tahmini için etkili araçlar olduğu gösterilmiştir.

Elde edilen sonuçlar:

\begin{itemize}
    \item Yıldırım tahmini için makine öğrenmesi yöntemlerinin başarıyla uygulanabileceğini
    \item Önceki günlerin yıldırım durumunun en önemli gösterge olduğunu
    \item Mevsimsel ve coğrafi faktörlerin tahmin gücüne katkı sağladığını
\end{itemize}

ortaya koymaktadır.

Bu çalışma, bölgesel yıldırım erken uyarı sistemleri için temel oluşturabilecek potansiyele sahiptir.
